\documentclass[11pt,a4paper]{article}

% ============================================================
% PACKAGES
% ============================================================
\usepackage[utf8]{inputenc}
\usepackage[T1]{fontenc}
\usepackage[english]{babel}
\usepackage[a4paper,margin=2.5cm]{geometry}
\usepackage{lmodern}
\usepackage{microtype}
\usepackage{amsmath,amssymb}
\usepackage{graphicx}
\usepackage{xcolor}
\usepackage{booktabs}
\usepackage{tabularx}
\usepackage{listings}
\usepackage{enumitem}
\usepackage{placeins}
\usepackage[hidelinks]{hyperref}

\setlength{\parindent}{0pt}
\setlength{\parskip}{0.5em}
\setlist[itemize]{nosep}
\setlist[enumerate]{nosep}

% ============================================================
% PLACEHOLDERS AND CODE STYLE
% ============================================================
\definecolor{codebackground}{RGB}{247,247,247}
\definecolor{codegreen}{RGB}{35,110,65}
\definecolor{codeblue}{RGB}{35,70,145}

\lstdefinestyle{CStyle}{
  language=C,
  basicstyle=\ttfamily\footnotesize,
  keywordstyle=\color{codeblue}\bfseries,
  commentstyle=\color{codegreen},
  stringstyle=\color{codegreen},
  backgroundcolor=\color{codebackground},
  numbers=left,
  numberstyle=\tiny\color{gray},
  numbersep=8pt,
  frame=single,
  rulecolor=\color{gray!35},
  breaklines=true,
  showstringspaces=false,
  keepspaces=true,
  columns=fullflexible,
  tabsize=4
}
\lstset{style=CStyle}

% Inserts an uploaded figure, or a placeholder if it is missing.
% Usage:
% \reportfigure{figures/example.png}{Caption}{fig:label}
\newcommand{\reportfigure}[3]{%
  \begin{figure}[htbp]
    \centering
    \IfFileExists{#1}{%
      \includegraphics[
        width=\linewidth,
        height=0.5\textheight,
        keepaspectratio
      ]{#1}%
    }{%
      \fbox{%
        \parbox[c][3.5cm][c]{0.90\linewidth}{%
          \centering
          \textbf{Figure placeholder}\\[0.5em]
          Upload the corresponding figure as\\[0.5em]
          \texttt{\detokenize{#1}}
        }%
      }%
    }
    \caption{#2}
    \label{#3}
  \end{figure}
}

% ============================================================
% DOCUMENT INFORMATION
% ============================================================
\title{
  \textbf{CPDS Laboratory Assignment 2}\\[0.4em]
  Parallel Programming with OpenMP\\
  \large Solving the Heat Equation
}
\author{
  Group member names not provided
}
\date{Fall 2026}

\begin{document}

\maketitle

\begin{center}
  \textbf{Group:} Not recorded\\
  \textbf{Submission date:} Not recorded
\end{center}

\tableofcontents
\newpage

% ============================================================
\section{Objectives and experimental setup}
% ============================================================

This report studies the parallelization of the Jacobi and
Gauss--Seidel heat diffusion solvers using OpenMP. We evaluate
correctness, strong scaling, synchronization overhead, and
memory locality for the implemented versions.

\subsection{Execution environment}

\begin{table}[htbp]
  \centering
  \begin{tabularx}{\linewidth}{lX}
    \toprule
    Item & Configuration \\
    \midrule
    Machine / execution node & \texttt{cpds1101} (recorded in Extrae table captions); processor model not recorded \\
    Available cores & 1, 4, 8, and 16 workers were measured; physical-core and SMT layout not recorded \\
    Compiler & Not recorded in the retained auxiliary files \\
    Compilation flags & Not recorded in the retained auxiliary files \\
    Thread placement & Not recorded in the retained auxiliary files \\
    Input configuration & Two heat sources; input filename and any changes not recorded \\
    Problem size & \texttt{-s 1022} in the user-parameter sweep \\
    Maximum iterations & \texttt{-n 1000} in the user-parameter sweep and 1000 solver invocations in traces \\
    Residual threshold & \texttt{-r 0.00005} in the user-parameter sweep \\
    Timing method & Extrae whole-program elapsed times for trace runs; ordinary elapsed-time output for the parameter sweep; repetitions not recorded \\
    \bottomrule
  \end{tabularx}
  \caption{Experimental facts recoverable from the retained auxiliary
  files; unavailable configuration details are identified in the table.}
  \label{tab:environment}
\end{table}

\subsection{Performance metrics}

Let $T_{\mathrm{seq}}$ be the execution time of the original
sequential implementation of the same solver, and let $T_p$
be the time of the parallel version using $p$ threads.
We define
\[
  S(p)=\frac{T_{\mathrm{seq}}}{T_p},
  \qquad
  E(p)=\frac{S(p)}{p}\times 100\%.
\]

If a plot instead uses the one-thread parallel execution as its
baseline, we label it explicitly:
\[
  S_1(p)=\frac{T_1}{T_p}.
\]

For comparisons between two implementations at the same thread
count, we use
\[
  G(p)=\frac{T_{\mathrm{before},p}}
             {T_{\mathrm{after},p}}.
\]

The retained Extrae model-factor outputs report whole-program elapsed
times for 1, 4, 8, and 16 workers. Their speedup and efficiency use the
one-worker run of the same version as the baseline, not a separately
measured sequential executable. The parameter sweep reports ordinary
runtime for 16 workers. The two timing sources are kept distinct below.
The sweep reached its 1000-iteration limit with residual 0.023071,
above the requested threshold 0.00005; it therefore did not converge
to that threshold.

\subsection{Implementation identifiers}

The identifiers below follow the versioned source and trace directories
in \texttt{aux/}. The source files combine Jacobi with one of the
Gauss--Seidel implementations.

\begin{table}[htbp]
  \centering
  \small
  \begin{tabularx}{\linewidth}{lXX}
    \toprule
    ID & Description & Source files \\
    \midrule
    J0 & Initial parallel Jacobi & v1 source \\
    J1 & Optimized parallel Jacobi & v2 source \\
    J2 & Jacobi with alternating input/output pointers & v3 source \\
    GT & Gauss--Seidel with explicit tasks & v1 source \\
    GD & Gauss--Seidel with do-across & v2 source \\
    GD-opt & Do-across with configurable column blocks & v3 source \\
    \bottomrule
  \end{tabularx}
  \caption{Jacobi and Gauss--Seidel variants identified by the auxiliary
  source-version directories.}
  \label{tab:versions}
\end{table}

Source basenames, in v1/v2/v3 order:\\
\url{solver-jacobi-parallel-gauss-explicit.c}\\
\url{solver-jacobi-optimized-gauss-do-across.c}\\
\url{solver-omp.c}

% ============================================================
\section{Sequential reference and correctness}
% ============================================================

\subsection{Sequential Jacobi and Gauss--Seidel}

Jacobi reads each iteration from one matrix and writes the new interior
values to a separate matrix, so every update uses the previous iteration's
neighbour values. The boundary cells are not updated. Gauss--Seidel
updates the matrix in place in row-major order; later cells in the sweep
can therefore observe values written earlier in that same iteration.
The driver stops when the residual is below the requested threshold or
when the maximum iteration count is reached (\texttt{-n 0} disables the
iteration limit in the supplied driver).

\begin{table}[htbp]
  \centering
  \begin{tabular}{lrrr}
    \toprule
    Solver & Time (s) & Iterations & Final residual \\
    \midrule
    Jacobi & Not recorded & Not recorded & Not recorded \\
    Gauss--Seidel & Not recorded & Not recorded & Not recorded \\
    \bottomrule
  \end{tabular}
  \caption{Sequential reference results for the stated configuration.}
  \label{tab:sequential}
\end{table}

No sequential-reference timings, iteration/residual logs, or PPM files
are present in the retained auxiliary files, so the sequential results
and output-image comparison cannot be reconstructed here.

\subsection{Validation procedure}

Each parallel implementation is compared with the original
sequential implementation of the \emph{same solver}, using
matching input and stopping conditions.

The trace and sweep artifacts do not contain paired sequential/parallel
PPM outputs or a correctness log. Consequently, the numerical equivalence
of each parallel version remains unverified by the retained evidence;
the model-factor traces demonstrate execution, not correctness.

\begin{table}[htbp]
  \centering
  \small
  \begin{tabularx}{\linewidth}{llrrX}
    \toprule
    Version & Threads & Iterations & Residual & PPM comparison \\
    \midrule
    J0 & 1, 8 & Not recorded & Not recorded & No PPM or comparison log retained \\
    J1 & 16 & Not recorded & Not recorded & No PPM or comparison log retained \\
    J2 & 1, 4, 8, 16 & 1000 invocations in trace & Not recorded & No PPM or comparison log retained \\
    GT & 1, 4, 8, 16 & 1000 invocations in trace & Not recorded & No PPM or comparison log retained \\
    GD & 1, 4, 8, 16 & 1000 invocations in trace & Not recorded & No PPM or comparison log retained \\
    GD-opt & 1, 4, 8, 16 & 1000 invocations in trace & Not recorded & No PPM or comparison log retained \\
    \bottomrule
  \end{tabularx}
  \caption{Configurations listed for validation in the draft. No
  pass/fail or PPM-comparison results were retained.}
  \label{tab:validation}
\end{table}

\FloatBarrier

% ============================================================
\section{Dependency analysis with Tareador}
% ============================================================

\subsection{Decompositions explored}

Tareador graph PDFs are retained for the Jacobi and Gauss--Seidel
examples under \texttt{tdg/}. The filenames distinguish the vanilla and
modified graph; the graph-generation command, exact problem size, and
simulation processor count were not retained. The modified Jacobi graph
shows potential tasks inside a solver invocation, while the vanilla graph
also shows the separate \texttt{copy\_mat} call.

The residual \texttt{sum} is a reduction in the OpenMP code. Any
dependencies deliberately omitted while exposing independent work to
Tareador must not be omitted from the real solver: the actual reduction
still needs a valid reduction or equivalent synchronization, and every
matrix dependency that protects a read-after-write must be preserved.

\subsection{Jacobi}

Within a Jacobi solve, all cells read only the old matrix and write
distinct interior cells in the new matrix, so cells can be computed
independently; only the residual reduction is shared. A copy from the
new matrix back to the old one (if used) must complete before the next
solve, and each iteration must complete before the following iteration
reads its result. These phase boundaries and repeated parallel-region
setup are serial/ synchronization costs that reduce useful parallel
work. The retained Tareador graph labels the solver and copy separately.

\paragraph{Tareador artifacts.}
The modified Jacobi graph is retained as
\url{tdg/dependency_graph_0_modified.pdf}; the vanilla graph, which
also shows the separate matrix copy, is
\url{tdg/dependency_graph_0_vanilla.pdf}. Their unusually large PDF
page dimensions make them unsuitable for embedding at readable scale in
this report.

The modified graph exposes multiple potential tasks per solver call,
but the saved graph is not a processor simulation and cannot establish
whether a particular processor count was sufficient. The critical path
must include the solver-to-copy-to-next-solver ordering; no simulation
timeline or processor-count metadata was retained.

\subsection{Gauss--Seidel}

For in-place Gauss--Seidel, a tile depends on the tile immediately to
its left and the tile immediately above it: those tiles produce values
read along the current tile's left and top boundaries. This imposes a
north-west to south-east wavefront; tiles on the same anti-diagonal can
run concurrently, while the first row and first column begin without one
of these predecessors.

\paragraph{Tareador artifacts.}
The modified Gauss--Seidel graph is retained as
\url{tdg/dependency_graph_1_modified.pdf}; the vanilla graph is
\url{tdg/dependency_graph_1_vanilla.pdf}. The original PDFs have page
dimensions too large for readable embedding in this report.

The Gauss--Seidel wavefront has less immediate parallelism near the
beginning and end of a sweep than Jacobi's independent cell updates.
The retained vanilla Gauss--Seidel graph has no potential task children,
whereas the modified graph exposes the block-level tasks. No simulation
artifact is available to quantify the resulting parallelism.

\FloatBarrier

% ============================================================
\section{Parallel Jacobi}
% ============================================================

\subsection{Initial data decomposition}

The \texttt{lowerb}, \texttt{upperb}, and \texttt{numElem} macros divide
the \texttt{sizex-2} interior rows among the threads, distributing any
remainder over the first blocks and offsetting the row indices by one.
Thus, the fixed boundary rows are excluded and each thread scans all
interior columns. The arrays are shared, but each thread writes a
disjoint set of \texttt{unew} cells. The scalar temporaries are private,
and \texttt{sum} is combined with an OpenMP reduction.



% The assignment requires explicit work assignment using
% OpenMP intrinsic functions for this Jacobi decomposition.
% Do not replace that requirement with omp for or explicit tasks.
\begin{lstlisting}[caption={Initial parallel Jacobi implementation.}]
 #pragma omp parallel private(tmp, diff) reduction(+:sum) 
        {

                int blocki=omp_get_thread_num();
                int nblocksi=omp_get_num_threads();
        
                int interior_rows = sizex - 2;
                int i_start = lowerb(blocki, nblocksi, interior_rows) + 1;
                int i_end = upperb(blocki, nblocksi, interior_rows) + 1;

                for (int i = i_start; i <= i_end; i++){
                        for (int j = 1; j<sizey-1; j++){
                                tmp = 0.25 * ( u[ i*sizey          + (j-1) ] +  // left
                                    u[ i*sizey     + (j+1) ] +  // right
                                    u[ (i-1)*sizey + j     ] +  // top
                                    u[ (i+1)*sizey + j     ] ); // bottom
                                diff = tmp - u[i*sizey+ j];
                                sum += diff * diff;
                                unew[i*sizey+j] = tmp;
          
        
                        }
                }       
                
                
        }       
\end{lstlisting}

No retained artifact records the requested 1- and 8-thread PPM or
residual comparison; see Table~\ref{tab:validation}.

\subsection{Overall analysis of the initial version}

The Extrae trace artifacts cover 1, 4, 8, and 16 workers and report
1000 solver invocations. The sweep independently records
\texttt{-s 1022}, \texttt{-n 1000}, and residual threshold 0.00005;
the exact trace command line is not present, so those input values
cannot be independently confirmed for every trace.

\begin{table}[htbp]
  \centering
  \begin{tabular}{rrrr}
    \toprule
    Workers & Time (s) & $S_1$ & $E_1$ (\%) \\
    \midrule
    1  & 2.46 & 1.00 & 100.0 \\
    4  & 1.62 & 1.51 & 37.8 \\
    8  & 1.47 & 1.67 & 20.9 \\
    16 & 1.50 & 1.63 & 10.2 \\
    \bottomrule
  \end{tabular}
  \caption{Initial Jacobi Extrae elapsed times and scalability relative
  to its own one-worker run.}
  \label{tab:jacobi-initial}
\end{table}

\reportfigure
  {aux/v1/heat-omp-jacobi-strong-extrae/modelfactor-tables.pdf}
  {Extrae model-factor tables for the v1 Jacobi trace.}
  {fig:jacobi-initial-mf}

The initial Jacobi trace reaches only 1.63x speedup at 16 workers and
its speedup peaks at 1.67x with 8 workers. The model reports a parallel
fraction of 73.83\%, global efficiency falling from 82.28\% at 4 workers
to 35.59\% at 16, and strategy efficiency near 98\% at 8 and 16 workers.
This points first to limited parallel fraction rather than poor task
balance: adding workers cannot accelerate the serial portion.

\subsection{Detailed analysis of the initial version}

The traces report 1000 implicit tasks per worker and no measured time in
implicit-task synchronization, but the retained directory has no
Paraver screenshots or synchronized timeline export. Thus the exact
source region responsible for the serial fraction cannot be identified
from the artifacts. The Tareador graph does show solver and copy phases
as separate operations, so iteration-boundary work is a plausible
candidate, not a confirmed attribution.

\subsection{Optimization}

\begin{lstlisting}[caption={Parallel interior-copy routine retained in the v1 source.}]
void copy_mat (double *u, double *v, unsigned sizex, unsigned sizey) {

    #pragma omp parallel
{

    int n_blocks = omp_get_num_threads();
    int blocki = omp_get_thread_num();

    int start_i = lowerb(blocki, n_blocks, sizex);
    int end_i = upperb(blocki, n_blocks,sizex); 

    for (int i = max(1, start_i); i <= min(sizex - 2, end_i) ; i++) {
        for (int j = 1; j < sizey - 1; j++) {
            v[i * sizey + j] = u[i * sizey + j];
        }
    }

}
}

\end{lstlisting}

The v2 source contains a serial \texttt{copy\_mat} routine instead.
The retained source files do not include the full build/run scripts,
so this parallel-copy code should not be assumed to be the exact v2
trace configuration.

No 16-worker correctness log or PPM comparison is retained.

\subsection{Overall analysis of the optimized version}

\begin{table}[htbp]
  \centering
  \begin{tabular}{rrrr}
    \toprule
    Workers & Time (s) & $S_1$ & $E_1$ (\%) \\
    \midrule
    1  & 2.44 & 1.00 & 100.0 \\
    4  & 0.68 & 3.58 & 89.5 \\
    8  & 0.36 & 6.74 & 84.3 \\
    16 & 0.20 & 12.38 & 77.4 \\
    \bottomrule
  \end{tabular}
  \caption{v2 Jacobi Extrae elapsed times and model-reported
  scalability relative to its own one-worker run.}
  \label{tab:jacobi-optimized}
\end{table}

\reportfigure
  {aux/v2/heat-omp-jacobi-strong-extrae/modelfactor-tables.pdf}
  {Model-factor tables for the optimized Jacobi implementation.}
  {fig:jacobi-optimized-mf}

For the v2 Jacobi trace, the 16-worker elapsed time is 0.20 s versus
1.50 s for v1, while the one-worker times are 2.44 s and 2.46 s.
The v2 trace's parallel fraction is 98.62\%, up from 73.83\% in v1.
Its trace speedup is 12.38x at 16 workers (versus 1.63x for v1).
The higher parallel fraction is supported by the model-factor output;
the retained files do not isolate which source-level change caused it.

The improvement over the initial implementation at 16 threads is
\[
  G_{\mathrm{Jacobi}}(16)
  =\frac{T_{\mathrm{J0},16}}{T_{\mathrm{J1},16}}
  =\frac{1.50}{0.20}=7.50.
\]

The v3 Jacobi trace reports 1.90, 0.49, 0.26, and 0.17 s for
1, 4, 8, and 16 workers respectively (speedups 1.00, 3.89, 7.28,
and 11.17). This is a separate implementation and is not the baseline
for the J1 comparison. No Jacobi trace shows superlinear speedup.

\subsection{Detailed analysis of the optimized version}

No Paraver timeline images or equivalent time-window markers are
retained, so a visual timeline comparison is not possible. The
model-factor tables show that v2's parallel fraction is much higher
than v1's, while global efficiency at 16 workers is still 81.65\%;
the remaining loss is reflected in the model's 84.89\% computation-task
scalability and 96.19\% strategy efficiency.

\subsection{Further optimization, if implemented}

The v3 source implements the pointer-alternation optimization: after a
Jacobi step, it swaps \texttt{param.u} and \texttt{param.uhelp} rather
than copying every interior element. The two buffers alternate roles;
the solver writes only interior cells, leaving each buffer's boundary
cells unchanged. The driver writes the buffer currently referenced by
\texttt{param.u} after the final iteration. This behavior is visible in
the source, but output-image correctness was not independently checked
in the retained artifacts.

\begin{lstlisting}[caption={Jacobi pointer alternation in the v3 driver.}]
case 0: // JACOBI
    residual = solve(param.u, param.uhelp, np, np);
    double *aux = param.u;
    param.u = param.uhelp;
    param.uhelp = aux;
    break;
\end{lstlisting}

The v3 trace reports 0.17 s with 16 workers, compared with 0.20 s
for v2, a 1.18x elapsed-time improvement (15\% lower). The retained
artifacts do not establish that the v2 and v3 runs used identical
commands or confirm correctness against a sequential PPM.

\FloatBarrier

% ============================================================
\section{Parallel Gauss--Seidel}
% ============================================================

\subsection{Blocking and dependency structure}

The Gauss--Seidel solver partitions the matrix into
\texttt{nblocksi} row blocks and \texttt{nblocksj} column blocks and
visits tiles in increasing row-block, then column-block, order.
The \texttt{lowerb}/\texttt{upperb} macros balance non-divisible
dimensions; each tile clips its iteration to the interior range
\texttt{1..sizex-2} by \texttt{1..sizey-2}. Since updates are in place,
each tile must wait for its left and upper neighbours. The resulting
anti-diagonal wavefront permits concurrent tiles without violating the
Gauss--Seidel read-after-write order.

\subsection{Task decomposition with explicit dependencies}

The explicit-task version creates one task per tile inside a
\texttt{single}/\texttt{taskgroup}. Each task owns a tile, keeps
\texttt{tmp} and \texttt{diff} private, and participates in the
taskgroup reduction for the residual. The two \texttt{depend(in:)}
locators represent the left and upper predecessor tiles; the
\texttt{depend(out:)} locator identifies the tile's bottom-right cell
for successor matching. The taskgroup completes all tile work and
combines the residual before the next solver iteration begins.

\begin{lstlisting}[caption={Explicit Gauss--Seidel tile tasks (v1).}]
int nblocksi = omp_get_max_threads();
int nblocksj = nblocksi;

#pragma omp parallel
{
  #pragma omp single
  {
    #pragma omp taskgroup task_reduction(+:sum)
    {
      for (int blocki = 0; blocki < nblocksi; ++blocki) {
        int i_start = lowerb(blocki, nblocksi, sizex);
        int i_end = upperb(blocki, nblocksi, sizex);
        for (int blockj = 0; blockj < nblocksj; ++blockj) {
          int j_start = lowerb(blockj, nblocksj, sizey);
          int j_end = upperb(blockj, nblocksj, sizey);
          #pragma omp task in_reduction(+:sum) private(tmp,diff) \
            depend(in: u[i_end*sizey+j_start-1], \
                       u[(i_start-1)*sizey+j_end]) \
            depend(out: u[i_end*sizey+j_end])
          {
            for (int i = max(1,i_start); i <= min(sizex-2,i_end); ++i)
              for (int j = max(1,j_start); j <= min(sizey-2,j_end); ++j) {
                tmp = 0.25 * (u[i*sizey+j-1] + u[i*sizey+j+1]
                              + u[(i-1)*sizey+j] + u[(i+1)*sizey+j]);
                diff = tmp - u[i*sizey+j];
                sum += diff * diff;
                u[i*sizey+j] = tmp;
              }
          }
        }
      }
    }
  }
}
\end{lstlisting}

The left locator matches the preceding tile's last-column output, and
the upper locator matches the preceding row tile's last-row output.
The output locator permits the next tile to depend on this tile.
However, source-code inspection is not a substitute for testing the
OpenMP dependency implementation; no sequential comparison or PPM
validation log is retained.

\begin{table}[htbp]
  \centering
  \begin{tabular}{rrrr}
    \toprule
    Workers & Time (s) & $S_1$ & $E_1$ (\%) \\
    \midrule
    1  & 7.59 & 1.00 & 100.0 \\
    4  & 4.24 & 1.79 & 44.8 \\
    8  & 2.85 & 2.66 & 33.2 \\
    16 & 2.01 & 3.78 & 23.6 \\
    \bottomrule
  \end{tabular}
  \caption{Explicit-task Gauss--Seidel Extrae elapsed times and
  scalability relative to its own one-worker run.}
  \label{tab:gauss-tasks}
\end{table}

\reportfigure
  {aux/v1/heat-omp-gausseidel-strong-extrae/modelfactor-tables.pdf}
  {Extrae model-factor tables for explicit-task Gauss--Seidel.}
  {fig:gauss-tasks-mf}

Explicit-task time falls from 7.59 s to 2.01 s, a 3.78x speedup at
16 workers. Load balance remains about 95\% for task count and task
execution time, but the model reports synchronization overhead per
task of 94.62\% at 4 workers and 153.88\% at 16; scheduling overhead
is only 1.68\% at 16. The workload creates 256,000 tasks at 16 workers
and 1000 taskgroups, so dependency/synchronization cost grows as each
tile's useful work shrinks. No Paraver timeline or memory-counter
results are retained to refine this interpretation.

\subsection{Data decomposition with do-across}

The do-across version distributes row-block iterations among the
implicit tasks and iterates over column blocks inside each row block.
The ordered source/sink clauses express the same left and upper
predecessor relations as the task version. Its reduction combines the
per-cell residual contributions.

\begin{lstlisting}[caption={Gauss--Seidel do-across dependency loop (v2).}]
int nblocksi = omp_get_max_threads();
int nblocksj = nblocksi;

#pragma omp parallel for ordered(2) private(tmp,diff) reduction(+:sum)
for (int blocki = 0; blocki < nblocksi; ++blocki) {
  int i_start = lowerb(blocki, nblocksi, sizex);
  int i_end = upperb(blocki, nblocksi, sizex);
  for (int blockj = 0; blockj < nblocksj; ++blockj) {
    int j_start = lowerb(blockj, nblocksj, sizey);
    int j_end = upperb(blockj, nblocksj, sizey);
    #pragma omp ordered depend(sink: blocki, blockj-1) \
                         depend(sink: blocki-1, blockj)
    for (int i = max(1,i_start); i <= min(sizex-2,i_end); ++i)
      for (int j = max(1,j_start); j <= min(sizey-2,j_end); ++j) {
        tmp = 0.25 * (u[i*sizey+j-1] + u[i*sizey+j+1]
                      + u[(i-1)*sizey+j] + u[(i+1)*sizey+j]);
        diff = tmp - u[i*sizey+j];
        sum += diff * diff;
        u[i*sizey+j] = tmp;
      }
    #pragma omp ordered depend(source)
  }
}
\end{lstlisting}

Each tile waits for its left and upper tile before updating in place,
then publishes completion as a source for its right and lower
successors. The v2 source uses an ordered loop nest over block indices;
the exact compiler/OpenMP version and a correctness run are not recorded,
so successful compilation and numerical equivalence cannot be confirmed
from the retained artifacts.

\begin{table}[htbp]
  \centering
  \begin{tabular}{rrrr}
    \toprule
    Workers & Time (s) & $S_1$ & $E_1$ (\%) \\
    \midrule
    1  & 7.58 & 1.00 & 100.0 \\
    4  & 3.29 & 2.30 & 57.5 \\
    8  & 1.74 & 4.36 & 54.5 \\
    16 & 0.90 & 8.43 & 52.7 \\
    \bottomrule
  \end{tabular}
  \caption{Initial do-across Gauss--Seidel Extrae elapsed times and
  scalability relative to its own one-worker run.}
  \label{tab:gauss-doacross}
\end{table}

\reportfigure
  {aux/v2/heat-omp-gausseidel-strong-extrae/modelfactor-tables.pdf}
  {Extrae model-factor tables for the initial do-across implementation.}
  {fig:gauss-doacross-mf}

At 16 workers, the do-across time is 0.90 s versus 2.01 s for explicit
tasks, while the 1-worker times (7.58 and 7.59 s) are nearly equal.
The do-across trace has 99.65\% parallel fraction and 99.37\% strategy
efficiency at 16 workers, but only 54.46\% computation-task scalability.
The recorded elapsed times show substantially better scaling than the
explicit-task version; no trace timeline is available to isolate where
the remaining loss occurs.

\FloatBarrier

% ============================================================
\section{Gauss--Seidel block-count tuning}
% ============================================================

\subsection{Configurable number of column blocks}

The assignment specifies
\[
  n_{\mathrm{blocks},j}
  = n_{\mathrm{blocks},i}\times \texttt{userparam}.
\]

In the v3 driver, \texttt{-u} sets the global \texttt{userparam}, which
the solver multiplies by the row-block count to obtain the number of
column blocks. The lower/upper-bound macros distribute remainder
elements across early blocks, while the tile loops clip their work to
interior cells.

\begin{lstlisting}[caption={Column-block count in the v3 solver.}]
int nblocksi = omp_get_max_threads();
int nblocksj = nblocksi * userparam;
\end{lstlisting}

Increasing \texttt{userparam} creates more tiles and exposes more
independent wavefront work, but reduces computation per tile and
increases dependency/scheduling frequency. Too many small tiles increase
overhead, as the 32- and 64-factor measurements demonstrate.

\subsection{Parameter exploration with 16 threads}

The retained sweep used 16 workers, resolution 1022, 1000 maximum
iterations, and residual threshold 0.00005. It reports ordinary
application elapsed time; the exact submission script and repetition
count are not present.

\begin{table}[htbp]
  \centering
  \begin{tabular}{rrrr}
    \toprule
    \texttt{userparam} &
    Row blocks &
    Column blocks &
    Time (s) \\
    \midrule
    1  & 16 & 16   & 0.922 \\
    2  & 16 & 32   & 0.634 \\
    4  & 16 & 64   & 0.480 \\
    8  & 16 & 128  & 0.421 \\
    16 & 16 & 256  & 0.464 \\
    32 & 16 & 512  & 0.665 \\
    64 & 16 & 1024 & 1.080 \\
    \bottomrule
  \end{tabular}
  \caption{Block-count exploration with 16 threads.
  All tested values from the retained sweep.}
  \label{tab:block-tuning}
\end{table}

Selected parameter:
\[
  \texttt{userparam}=8.
\]

Among the tested values, 8 is fastest at 0.421 s. Increasing it to 16,
32, and 64 worsens runtime to 0.464, 0.665, and 1.080 s, consistent
with synchronization and task overhead outweighing additional
parallelism. The sweep records the same 1000-iteration residual
(0.023071) for every value, but contains no PPM comparison or independent
correctness check.

\subsection{Overall analysis after tuning}

The retained v3 driver defaults \texttt{userparam} to 1; the trace files
do not preserve their command line, so it is not possible to confirm
that the selected value 8 was used for the strong-scaling traces.
With a fixed parameter of 8, column-block count would scale as eight
times the worker count (8, 32, 64, and 128 for 1, 4, 8, and 16 workers).

\begin{table}[htbp]
  \centering
  \begin{tabular}{rrrr}
    \toprule
    Workers & Time (s) & $S_1$ & $E_1$ (\%) \\
    \midrule
    1  & 7.30 & 1.00 & 100.0 \\
    4  & 1.65 & 4.44 & 111.0 \\
    8  & 0.71 & 10.23 & 127.9 \\
    16 & 0.40 & 18.24 & 114.0 \\
    \bottomrule
  \end{tabular}
  \caption{v3 do-across Gauss--Seidel Extrae elapsed times and reported
  scalability relative to its own one-worker run. The trace's
  \texttt{-u} value is unknown.}
  \label{tab:gauss-tuned}
\end{table}

\reportfigure
  {aux/v3/heat-omp-gausseidel-strong-extrae/modelfactor-tables.pdf}
  {Extrae model-factor tables for the v3 do-across implementation.}
  {fig:gauss-tuned-mf}

The v3 trace reports 18.24x speedup at 16 workers, above linear; its
model reports 121.72\% global efficiency and 123.57\% computation-task
scalability, driven in part by 157.12\% IPC scalability. Instruction
scalability is 84.22\% and frequency scalability 93.39\%, so this
superlinear result should be treated as measurement/model behavior,
not as evidence of more than ideal parallel work. Compared with v2's
0.90 s at 16 workers, v3's 0.40 s is 2.25x faster, but the v3 trace's
\texttt{-u} value is not recorded and comparability is therefore
conditional.

\FloatBarrier

% ============================================================
\section{Memory hierarchy analysis}
\label{sec:memory}
% ============================================================

\subsection{Measurement configuration}

Raw Extrae trace files are retained for the explicit-task and
do-across versions on \texttt{cpds1101}, but the memory-measurement
script, counter configuration, and derived L2-miss output are absent.
The traces alone do not establish which cache event was measured or
whether the parameter and iteration counts match the timing runs.

\subsection{L2 cache misses}

\begin{table}[htbp]
  \centering
  \small
  \begin{tabular}{lrrrr}
    \toprule
    Version & Threads & Total L2 misses &
    Misses/thread & Time (s) \\
    \midrule
    \multicolumn{5}{c}{No L2-miss counts or matching memory timings retained} \\
    \bottomrule
  \end{tabular}
  \caption{No cache-counter or matching timing values were retained
  for explicit tasks or data decomposition.}
  \label{tab:memory}
\end{table}

The retained artifacts do not report total or per-thread L2 misses, nor
do they contain the matching memory-run completion counts. The parameter
sweep reaches 1000 iterations with residual 0.023071, but it is not an
L2-counter measurement and cannot stand in for one.

\subsection{Interpretation}

No quantitative cache-locality conclusion is supported without the
missing L2 counter values. By construction, row ownership gives the
data-decomposition versions contiguous row segments, while explicit
task scheduling adds task/dependency overhead; these are hypotheses
about locality and overhead, not measured cache results here.

\FloatBarrier

% ============================================================
\section{Conclusions}
% ============================================================

Among the retained Jacobi traces, the v3 pointer-alternation version
has the lowest 16-worker elapsed time (0.17 s), compared with 0.20 s
for v2 and 1.50 s for v1. The model's reported parallel fraction rises
from 73.83\% in v1 to 98.62\% in v2. For Gauss--Seidel, the explicit
task version reaches 3.78x speedup at 16 workers, while the initial
do-across trace reaches 8.43x; the v3 trace reports 18.24x, subject to
the unknown \texttt{-u} value and the superlinear-efficiency caveat.
The best parameter in the retained 16-worker sweep is
\texttt{userparam}=8 at 0.421 s. No cache-counter evidence is available.

The reported scaling depends on the traced problem size and thread
count, and the available records do not specify processor topology,
compiler flags, repetitions, or the exact command line for every trace.
The sweep failed to reach its requested residual threshold within 1000
iterations. Correctness comparisons, Paraver timelines, and memory
counters should be added before treating these performance observations
as a complete evaluation.

% Submission reminder:
% Submit the report PDF and a separate ZIP containing only
% the parallel C source files, with meaningful version names.
% Include every group member's name on the title page.
% Replace the unrecorded group metadata before submission.

\end{document}
